\section{Related Work}
\paragraph{Flow metrics and developer productivity}
\ffrs{}'s metrics are descendants of established flow vocabularies. The DevOps programme summarised in \emph{Accelerate}~\cite{forsgren2018accelerate} and later expanded in the SPACE framework~\cite{forsgren2021space} popularised outcome metrics---lead time for changes, deployment frequency---chosen because they can be measured mechanically. ITIL~4~\cite{axelos2019itil} frames request handling as a sequence of acknowledgement, routing, and resolution targets. \ffrs{} borrows the shape of both but reduces the scope to what one tracker can carry: no service desk, no separate measurement system. Where DORA measures the path from commit to production, \ffrs{} measures the path from stakeholder to outcome; where ITIL assumes a staffed queue, \ffrs{} assumes an agent produces the first response.

\paragraph{Issue handling and maintainer fatigue}
Large-scale mining studies established both the centrality and fragility of issue trackers. Bissyand\'e et al.~\cite{bissyande2013gotissues} found that responsiveness varies widely, while Kalliamvakou et al.~\cite{kalliamvakou2014promises} catalogued the pitfalls of interpreting tracker data. Unmanaged queues are among the reasons open source projects fail~\cite{coelho2017why}. \ffrs{} takes these findings as motivation: the tracker is where decisions are made, so \ffrs{} keeps it as the system of record and instruments it, assigning the first response to an agent to mitigate fatigue.

\paragraph{Autonomous issue resolution}
SWE-bench~\cite{jimenez2024swebench} made ``resolve a real GitHub issue with a pull request'' a benchmark task. SWE-agent~\cite{yang2024sweagent}, Agentless~\cite{xia2024agentless}, ChatDev~\cite{qian2023chatdev}, and RepairAgent~\cite{bouzenia2024repairagent} have demonstrated that language models with purpose-built interfaces can resolve issues unattended at meaningful rates. \ffrs{} does not advance agent capability. Instead, it positions the agent as the \emph{default} owner of the Respond stage in a feedback loop, bounded by human checkpoints (merge for code; requester acceptance and reviewer confirmation for proposals), and it measures the result with the same timestamps as human-only handling. To our knowledge the combination---a feedback pipeline model, an agentic first response, and metrics recomputable from the public tracker---has not been described as a unit.

\paragraph{In-product feedback tooling and mining}
Extracting features and bugs from user feedback is a rich area of research, with tools like SAFE~\cite{johann2017safe} and various machine learning classifiers for app reviews~\cite{panichella2015how, guzman2014how}. While commercial widgets (e.g., Sentry, Intercom) provide the capture surface, they rely on vendor-hosted storage. \ffrs{} sits between these: an anonymous, account-free capture surface that funnels directly into the public tracker, avoiding proprietary databases while keeping personal data in a private sidecar.
